\vspace*{\fill}
\begin{glyph}{Mountain (山)}{mountain}\label{glyph:mountain}
Mountain.\\\hline
Obvious from the glyph at first glance.
\end{glyph}

\begin{comps}
\comprtwocone & 山 \\\hline
\comprthreecone & Mountain \\\hline
\end{comps}

\begin{remglyph}
Part of the Kanji radicals (\#{46})\\\hline
\end{remglyph}

\begin{prons}
\pronsrtwocone & \textbf{サン (SAN)} & \textbf{やま (yama)} \\\hline
\pronsrthreecone & --- & --- \\\hline
\end{prons}

\begin{remprons}
On the \hooglig{mountain}, my \comon{son} played a \comkun{yamaha} piano.\\\hline
\end{remprons}

\begin{somme}
    山 & \textit{mountain} & やま \mbox{(yama)} \\\hline
    \notyet{一山} & one + mountain = \textit{a pile of something} &  ひと{\middel}やま \mbox{(hito{\middel}yama)}\\\hline
    \notyet{高山} & tall + mountain = \textit{alpine} & コウ{\middel}ザン \mbox{(K\={O}{\middel}ZAN)}\\\hline
    \notyet{火山} & fire + mountain = \textit{volcano} & カ{\middel}ザン \mbox{(KA{\middel}ZAN)}\\\hline
    \notyet{山村} & mountain + village = \textit{mountain village} & サン{\middel}ソン \mbox{(SAN{\middel}SON)}\\\hline
\end{somme}

\begin{sentence}
\ruby{\underline{山}}{やま}に\ruby{\underline{木}}{き}が\underline{ありません}。\\\hline
yama ni ki ga arimasen.\\\hline
(mountain) (tree) (are not)\\\hline
=\textit{There are no trees on the mountain.}\\\hline
\end{sentence}
\end{multicols}
\vhrulefill{5pt}
%--------------------------------------------------------------------
%--------------------------------------------------------------------
\begin{multicols}{2}
\begin{glyph}{Person (人)}{person}\label{glyph:person}
Person; Human being. \\\hline
When appearing at the end of country names, this kanji denotes an individual's nationality.
\end{glyph}

\begin{comps}
\comprtwocone & 人 \\\hline
\comprthreecone & Person \\\hline
\end{comps}

\begin{remglyph}
Part of the Kanji radicals (\#{9}).\\\hline
\end{remglyph}

\begin{prons}
\pronsrtwocone & \textbf{ジン、ニン (JIN, NIN)} & \textbf{ひと (hito)} \\\hline
\rye{3}{ジン tends to signify that a person belongs to a certain subgroup of
	humanity.  Look for ジン in the first position and at the end of words signifying countries.\\\hline
	ニン tends to indicate a person engaged in an activity specified by the kanji preceding it.\\\hline
	The 訓読み becomes always voiced when it appears outside of the first position.}
\pronsrthreecone & --- & --- \\\hline
\end{prons}

\begin{remprons}
That \hooglig{person} is a \comon{nincompoop} because he \comon{jinxes} with
    \comkun{hit-or-miss} attempts.\\\hline
\end{remprons}

\begin{somme}
    人 & \textit{person} & ひと \mbox{(hito)} \\\hline
    \notyet{白人} & white + person = \textit{caucasian} & ハク{\middel}ジン \mbox{(HAKU{\middel}JIN)} \\\hline
    \notyet{人口} & person + mouth = \textit{population} & ジン{\middel}コウ \mbox{(JIN{\middel}K\={O})} \\\hline
    \notyet{三人} & three + person = \textit{three people} & サン{\middel}ニン \mbox{(SAN{\middel{NIN}})}\\\hline
    \notyet{四人} & four + person = \textit{four people} & よ{\middel}{ニン} \mbox{(yo{\middel}{NIN})}\\\hline
    \notyet{村人} & village + person = \textit{villager} & むら{\middel}{びと} \mbox{(mura{\middel}bito)}\\\hline
    \notyet{外国人} & outside + country + person = \textit{foreigner} & ガイ{\middel}コク{\middel}ジン \mbox{(GAI{\middel}KOKU{\middel}JIN)}\\\hline
\end{somme}

\begin{irrr}
\rye{3}{The irregular readings in the first two examples belong to 人.  In the final compound, both readings are irregular.}
    \notyet{一人} & one + person = \textit{one person} & ひと{\middel}り \mbox{(hito{\middel}ri)}\\\hline
    \notyet{二人} & two + person = \textit{two people} & ふた{\middel}り \mbox{(futa{\middel}ri)}\\\hline
    \notyet{大人} & large + person = \textit{adult} & おとな \mbox{(otona)}\\\hline
\end{irrr}

\begin{sentence}
\underline{あの} \ruby{\underline{人}}{ひと} の\ruby{\underline{手}}{て}\ruby{は}{わ}\ruby{\underline{小}}{ちい}\underline{さい}です。\\\hline
ano hito no te wa chii{\middel}sai desu.\\\hline
(that) (person) (hand) (small).\\\hline
=\textit{That person's hands are small.}\\\hline
\end{sentence}
\end{multicols}
\vhrulefill{5pt}
%--------------------------------------------------------------------
%--------------------------------------------------------------------
\begin{multicols}{2}
\begin{glyph}{One (一)}{one}\label{glyph:one}
One.
\end{glyph}

\begin{comps}
\comprtwocone & 一 \\\hline
\comprthreecone & One \\\hline
\end{comps}

\begin{remglyph}
Part of the Kanji radicals (\#{1}).\\\hline
\end{remglyph}

\begin{prons}
\pronsrtwocone & \textbf{イチ、イツ (ICHI, ITSU)} & \textbf{ひと (hito)} \\\hline
\pronsrthreecone & --- & --- \\\hline
\rye{3}{When this kanji appears as the first character in a compound it is almost always read with its on-yomi.  \\\hline
    Its pronounciation イチ or イツ will depend on whether the initial sound of the following kanji is voiced or voiceless.}
\end{prons}

\begin{remprons}
\hooglig{One} \comkun{hit-or-miss} attempt to \comon{eat soup} that reduces \comon{itchiness}.\\\hline
\end{remprons}

\begin{somme}
    一 & \textit{one} & イチ \mbox{(ICHI)}\\\hline
    一つ & \textit{one (general counter)} & ひと{\middel}つ \mbox{(hito{\middel}tsu)}\\\hline
    \notyet{一日} & one + sun (day) = \textit{one day} & イチ{\middel}ニチ \mbox{(ICHI{\middel}NICHI)}\\\hline
    \notyet{一月} & one + moon (month) = \textit{January} & イチ{\middel}ガツ \mbox{(ICHI{\middel}GATSU)}\\\hline
    \notyet{同一} & same + one = \textit{identical} & ドウ{\middel}イツ \mbox{(D\={O}{\middel}ITSU)}\\\hline
    \notyet{一回} & one + rotate = \textit{once} & イッ{\middel}カイ \mbox{(IK{\middel}KAI)}\\\hline
    \notyet{一時} & one + time = \textit{one o'clock} & イチ{\middel}ジ \mbox{(ICHI{\middel}JI)}\\\hline
\end{somme}

\begin{irrr}
    \notyet{一日} & one + sun (day) = \textit{first day of the month} & ついたち (tsuitachi)\\\hline
\rye{3}{The same compound is shown above with a different pronunciation and meaning.  Context will determine which is the appropriate reading to apply.  For example, when 一日 is seen at the top of a newspaper page, it clearly refers to the first day of the month.}
\end{irrr}

\begin{sentence}
\ruby{\underline{一時}}{イチジ}に\ruby{\underline{会}}{あ}\underline{いましょう}。\\\hline
ICHI{\middel}JI ni a{\middel}imash\={o}.\\\hline
(one o'clock) (let's meet)\\\hline
=\textit{Let's meet at one o'clock.}\\\hline
\end{sentence}
\end{multicols}
\vspace*{\fill}
%%--------------------------------------------------------------------
%%--------------------------------------------------------------------
\pagebreak
\begin{multicols}{2}
\begin{glyph}{Two (二)}{two}\label{glyph:two}
Two.  Including the ideas of ``double'' and ``bi-'', etc.
\end{glyph}

\begin{comps}
\comprtwocone & 一 + 一 \\\hline
\comprthreecone & One + One \\\hline
\end{comps}

\begin{remglyph}
Part of the Kanji radicals (\#{7}).\\\hline
\end{remglyph}

\begin{prons}
\pronsrtwocone & \textbf{ニ (NI)} & \textbf{ふた (futa)} \\\hline
\pronsrthreecone & --- & --- \\\hline
\end{prons}

\begin{remprons}
    \hooglig{Two} \comon{nickels} to tell you \comkun{who tans}?\\\hline
\end{remprons}

\begin{somme}
    ニ & \textit{two} & ニ \mbox{(NI)}\\\hline
    二つ & \textit{two (general counter)} & ふた{\middel}つ \mbox{(futa{\middel}tsu)}\\\hline
    \notyet{二月} & two + moon (month) = \textit{February} & ニガ{\middel}ツ \mbox{(NI{\middel}GATSU)}\\\hline
    \notyet{二十} & two + ten = \textit{twenty} & ニ{\middel}ジュウ \mbox{(NI{\middel}J\={U})}\\\hline
    \notyet{二百} & two + hundred = \textit{two hundred} & ニ{\middel}ヒャク \mbox{(NI{\middel}HYAKU)}\\\hline
    \notyet{二時} & two + time = \textit{two o'clock} & ニ{\middel}ジ \mbox{(NI{\middel}JI)}\\\hline
    \notyet{二週間} & two + week + interval = \textit{two weeks} & ニ{\middel}シュウ{\middel}カン \mbox{(NI{\middel}SH\={U}{\middel}KAN)}\\\hline
\end{somme}

\begin{irrr}
    \notyet{二日} & two + sun (day) = \textit{second day of the month} & ふつか \mbox{(futsuka)}\\\hline
    \notyet{二十日} & two + ten + sun (day) = \textit{twentieth day of the month} & はつか \mbox{(hatsuka)}\\\hline
    \notyet{二十歳} & two + ten + annual = \textit{twenty years old} & はたち \mbox{(hatachi)}\\\hline
\end{irrr}

\begin{sentence}
\ruby{\underline{二月}}{ニガツ}の\underline{オーストラリア}\ruby{は}{わ}\ruby{\underline{美}}{うつく}\underline{しい}。\\\hline
NI{\middel}GATSU no \={o}sutoraria wa utsuku{\middel}shii.\\\hline
(February) (Australia) (beautiful)\\\hline
=\textit{Australia is beautiful in February.}\\\hline
\end{sentence}
\end{multicols}
\vhrulefill{5pt}
%%--------------------------------------------------------------------
%%--------------------------------------------------------------------
\begin{multicols}{2}
\begin{glyph}{Three (三)}{three}\label{glyph:three}
Three/Triple/Tri-, etc.
\end{glyph}

\begin{comps}
\comprtwocone & 一 + 一 + 一 \\\hline
\comprthreecone & One + One + One \\\hline
\end{comps}

\begin{remglyph}
---\\\hline
\end{remglyph}

\begin{prons}
\pronsrtwocone & \textbf{サン (SAN)} & --- \\\hline
\pronsrthreecone & --- & みっ、み (mi-, mi) \\\hline
\end{prons}

\begin{remprons}
My \hooglig{triplet} \comon{sons} \ucomkun{mirror} one another.\\\hline
\end{remprons}

\begin{somme}
    三 & \textit{three} & サン \mbox{(SAN)}\\\hline
    三つ & \textit{three (general counter)} & みっ{\middel}つ \mbox{(mit{\middel}tsu)}\\\hline
    \notyet{三人} & three + person = \textit{three people} & サン{\middel}ニン \mbox{(SAN{\middel}NIN)}\\\hline
    \notyet{三月} & three + moon (month) = \textit{March} & サン{\middel}ガツ \mbox{(SAN{\middel}GATSU)}\\\hline
    \notyet{三時} & three + time = \textit{three o'clock} & サン{\middel}ジ \mbox{(SAN{\middel}JI)}\\\hline
    \notyet{三十} & three + ten = \textit{thirty} & サン{\middel}ジュウ \mbox{(SAN{\middel}J\={U})}\\\hline
    \notyet{三千} & three + thousand = \textit{three thousand} & サン{\middel}ゼン \mbox{(SAN{\middel}ZEN)}\\\hline
    \notyet{三日} & three + sun (day) = \textit{third day of the month} & みっ{\middel}か \mbox{(mik{\middel}ka)}\\\hline
\end{somme}

\begin{sentence}
\ruby{\underline{三人}}{サンニン}が\ruby{\underline{早}}{はや}\underline{く} \ruby{\underline{来}}{き}\underline{ました}。\\\hline
SAN{\middel}NIN ga haya{\middel}ku ki{\middel}mashita.\\\hline
(three people) (early) (came).\\\hline
=\textit{Three people came early.}\\\hline
\end{sentence}
\end{multicols}
%%--------------------------------------------------------------------
%%--------------------------------------------------------------------
\pagebreak
\begin{multicols}{2}
\begin{glyph}{Sun (日)}{sun}\label{glyph:sun}
Sun, and Day.\\\hline
It is the initial kanji in the compound for Japan.
\end{glyph}

\begin{comps}
\comprtwocone & 日 \\\hline
\comprthreecone & Sun \\\hline
\end{comps}

\begin{remglyph}
Part of the Kanji radicals (\#{72})\\\hline
\end{remglyph}

\begin{prons}
\pronsrtwocone & \textbf{ニチ、ジツ (NICHI, JITSU)} & \textbf{ひ (hi)} \\\hline
\pronsrthreecone & --- & か (ka) \\\hline
\rye{3}{ニチ will be encountered overwhelmingly in the first position (ひ only makes a few appearances).\\\hline
ジツ is typically encountered in the second position.\\\hline
ひ is typically encountered in the third position.\\\hline
か is used only to identify the days of the month from 2--10, 14, 17, 20, 24, and 27.}
\end{prons}

\begin{remprons}
The \hooglig{sun} inspired \comon{Nichiren} \comon{Ju-Jitsu}; many other
    \ucomkun{customs} are \comkun{heaped} under the \hooglig{sun}.\\\hline
\end{remprons}

\begin{somme}
    日 & \textit{sun; day} & ひ \mbox{(hi)}\\\hline
    一日 & one + day = \textit{one day} & イチ{\middel}ニチ \mbox{(ICHI{\middel}NICHI)} \\\hline
    \notyet{日本} & sun + main = \textit{Japan} & ニ{\middel}ホン、ニッ{\middel}ポン \mbox{(NI{\middel}HON, NIP{\middel}PON)}\\\hline
    \notyet{休日} & rest + day = \textit{holiday} & キュウ{\middel}ジツ \mbox{(KY\={U}{\middel}JITSU)}\\\hline
    \notyet{毎日} & every + day = \textit{daily} & マイ{\middel}ニチ \mbox{(MAI{\middel}NICHI)}\\\hline
    \notyet{朝日} & morning + sun = \textit{morning sun} & あさ{\middel}ひ \mbox{(asa{\middel}hi)}\\\hline
    \notyet{日曜日} & sun + day of the week + day = \textit{Sunday} & ニチ{\middel}ヨウ{\middel}ビ \mbox{(NICHI{\middel}Y\={O}{\middel}bi)}\\\hline
\end{somme}

\begin{irrr}
    一日 & one + day = \textit{first day of the month} & ついたち \mbox{(tsuitachi)}\\\hline
    \notyet{昨日} & past + day = \textit{yesterday} & きのう \mbox{(kin\={o})}\\\hline
    \notyet{今日} & now + day = \textit{today} & きょう \mbox{(ky\={o})}\\\hline
    \notyet{明日} & bright + day = \textit{tomorrow} & あした、あす \mbox{(ashita, asu)}\\\hline
\end{irrr}

\begin{sentence}
\ruby{\underline{毎日}}{マイニチ}{\;}\ruby{\underline{日本語}}{ニホンゴ}の\ruby{\underline{本}}{ホン}\ruby{を}{お}\ruby{\underline{読}}{よ}\underline{みます}。\\\hline
MAI{\middel}NICHI NI{\middel}HON{\middel}GO no HON o yo{\middel}mimasu.\\\hline
(daily) (Japanese) (book) (read).\\\hline
=\textit{(I) read Japanese books everyday.}\\\hline
\end{sentence}
\end{multicols}
\vhrulefill{5pt}
%%--------------------------------------------------------------------
%%--------------------------------------------------------------------
\begin{multicols}{2}
\begin{glyph}{White (白)}{white}\label{glyph:white}
White.  As well as shades of cleanliness and purity.
\end{glyph}

\begin{comps}
\comprtwocone & 白 \\\hline
\comprthreecone & White \\\hline
\end{comps}

\begin{remglyph}
    Part of the Kanji radicals (\#{106}).\\\hline
\end{remglyph}

\begin{prons}
\pronsrtwocone & \textbf{ハク (HAKU)} & \textbf{しろ (shiro)} \\\hline
\pronsrthreecone & ビャク (BYAKU) & しら (shira) \\\hline
\end{prons}

\begin{remprons}
    In \hooglig{white} robes, the \comkun{sheerly obstinate} \ucomon{Byakuya}
    challenged \comon{Haku}; it was \ucomkun{sheer ugliness}.\\\hline
\end{remprons}

\begin{somme}
    白 & \textit{white (noun)} & しろ \mbox{(shiro)}\\\hline
    白い & \textit{white (adjective)} & しろ{\middel}い \mbox{(shiro{\middel}i)}\\\hline
    白人 & white + person = \textit{caucasian} & ハク{\middel}ジン \mbox{(HAKU{\middel}JIN)}\\\hline
    \notyet{明白} & bright + white = \textit{obvious} & メイ{\middel}ハク \mbox{(MEI{\middel}HAKU)}\\\hline
    \notyet{白鳥} & white + bird = \textit{swan} & ハク{\middel}チョウ \mbox{(HAKU{\middel}CH\={O})}\\\hline
    \notyet{白米} & white + rice = \textit{polished rice} & ハク{\middel}マイ \mbox{(HAKU{\middel}MAI)}\\\hline
\end{somme}

\begin{sentence}
\underline{この} \ruby{\underline{白}}{しろ}\underline{い} \ruby{\underline{馬}}{うま}の\ruby{\underline{名前}}{なまえ}\ruby{は}{わ}\ruby{\underline{``雪''}}{ゆき}です。\\\hline
kono shiro{\middel}i uma no na{\middel}mae wa yuki desu.\\\hline
(this) (white) (horse) (name) (``snow'').\\\hline
=\textit{This white horse's name is ``Snow''}\\\hline
\end{sentence}
\end{multicols}
%%--------------------------------------------------------------------
%%--------------------------------------------------------------------
\pagebreak
\vspace*{\fill}
\begin{multicols}{2}
\begin{glyph}{Mouth (口)}{mouth}\label{glyph:mouth}
Mouth, and things oral.  This kanji also relates to openings in general, from cavers and harbors to taps and bottles, for instance.
\end{glyph}

\begin{comps}
\comprtwocone & 口 \\\hline
\comprthreecone & Mouth \\\hline
\end{comps}

\begin{remglyph}
Part of the Kanji radicals (\#{30}).\\\hline
\end{remglyph}

\begin{prons}
\pronsrtwocone & \textbf{コウ (K\={O})} & \textbf{くち (kuchi)} \\\hline
\pronsrthreecone & ク (KU) & --- \\\hline
\end{prons}

\begin{remprons}
His \hooglig{mouth} \comon{caused} the \ucomon{cook} to \comkun{coochi-coo} at
    him.\\\hline
\end{remprons}

\begin{somme}
    口 & \textit{mouth} & くち \mbox{(kuchi)}\\\hline
    人口 & person + mouth = \textit{population} & ジン{\middel}コウ \mbox{(JIN{\middel}K\={O})}\\\hline
    \notyet{早口} & early (fast) + mouth = \textit{rapid speaking} & はや{\middel}くち \mbox{(haya{\middel}kuchi)}\\\hline
    \notyet{口語} & mouth + words = \textit{colloquial language} & コウ{\middel}ゴ \mbox{(K\={O}{\middel}GO)}\\\hline
    \notyet{出口} & exit + mouth = \textit{exit} & で{\middel}ぐち \mbox{(de{\middel}guchi)}\\\hline
    \notyet{口先} & mouth + precede = \textit{lip service} & くち{\middel}さき \mbox{(kuchi{\middel}saki)}\\\hline
\end{somme}

\begin{sentence}
\underline{この} \ruby{\underline{国}}{くに}の\ruby{\underline{人口}}{ジンコウ}\ruby{は}{わ}\ruby{\underline{少}}{すく}\underline{ない}。\\\hline
kono kuni no JIN{\middel}K\={O} wa suku{\middel}nai.\\\hline
(this) (country) (population) (few).\\\hline
=\textit{This country has a small population.}\\\hline
\end{sentence}
\end{multicols}
\vhrulefill{5pt}
%%--------------------------------------------------------------------
%%--------------------------------------------------------------------
\begin{multicols}{2}
\begin{glyph}{Rotate (回)}{rotate}\label{glyph:rotate}
The general sense of this kanji relates to the ideas of rotation and things going around in both space and time.\\\hline
It can also refer to the vicinity of things (neighbourhoods, surroundings, etc.).\\\hline
This kanji is well-known to Japanese baseball fans as the character used to denote innings.
\end{glyph}

\begin{comps}
\comprtwocone & 囗 + 口 \\\hline
\comprthreecone & Enclosure + Mouth \\\hline
\end{comps}

\begin{remglyph}
\hooglig{Rotate} \remcom{dracula} in the \remcom{enclosure}.\\\hline
\end{remglyph}

\begin{prons}
\pronsrtwocone & \textbf{カイ (KAI)} & \textbf{まわ (mawa)} \\\hline
\pronsrthreecone & --- & --- \\\hline
\end{prons}

\begin{remprons}
The \hooglig{rotating} \comon{kites} that \comkun{ma wacks}.\\\hline
\end{remprons}

\begin{somme}
\rye{3}{Note the difference between intransitive (intr.) and transitive (tr.) verbs.  An intransitive verb does not require an object to complete the action, whereas a transitive verb does.}
    回る (intr.) & \textit{to go around} & まわ{\middel}る \mbox{(mawa{\middel}ru)}\\\hline
    回す (tr.) & \textit{to send around} & まわ{\middel}す \mbox{(mawa{\middel}su)}\\\hline
    一回 & one + rotate = \textit{once} & イッ{\middel}カイ \mbox{(IK{\middel}KAI)}\\\hline
    二回 & two + rotate = \textit{twice} & ニ{\middel}カイ \mbox{(NI{\middel}KAI)}\\\hline
    \notyet{引き回す} & pull + rotate = \textit{to pull around} & ひ{\middel}き　まわ{\middel}す \mbox{(hi{\middel}ki mawa{\middel}su)}\\\hline
    \notyet{言い回し} & say + rotate = \textit{turn of expression} & い{\middel}い　まわ{\middel}し \mbox{(i{\middel}i mawa{\middel}shi)}\\\hline
\end{somme}

\begin{sentence}
\ruby{\underline{私}}{わたし}\ruby{は}{わ}\ruby{\underline{二回}}{ニカイ}{\;}\underline{カナダ}\ruby{へ}{え}\ruby{\underline{行}}{い}\underline{きました}。\\\hline
watashi wa NI{\middel}KAI Canada e i{\middel}kimashita.\\\hline
(I) (twice) (Canada) (went)\\\hline
I went to Canada twice.\\\hline
\end{sentence}
\end{multicols}
\vspace*{\fill}
%%--------------------------------------------------------------------
%%--------------------------------------------------------------------
\pagebreak
\vspace*{\fill}
\begin{multicols}{2}
\begin{glyph}{Four (四)}{four}\label{glyph:four}
Four, Quad-, etc.
\end{glyph}

\begin{comps}
\comprtwocone & 囗 + 儿 \\\hline
    \comprthreecone & Enclosure + Person's legs
    (\ref{glyph:person}) \\\hline
\end{comps}

\begin{remglyph}
\hooglig{Four} \remcom{people} danced to the \hooglig{four} corners of the
    \remcom{prison cell}.\\\hline
\end{remglyph}

\begin{prons}
\pronsrtwocone & \textbf{シ (SHI)} & \textbf{よん、よ (yon, yo)} \\\hline
\pronsrthreecone & --- & よっ (yo-) \\\hline
\end{prons}

\begin{remprons}
\hooglig{Four} \comon{shields} \comkun{yomped} \comkun{yonder} to the
    \hooglig{four} corners of the world.\\\hline
\end{remprons}

\begin{somme}
    四 & \textit{four} & シ \mbox{(SHI)} \\\hline
    四人 & four + person = \textit{four people} & よ{\middel}ニン \mbox{(yo{\middel}NIN)} \\\hline
    四日 & four + day = \textit{the fourth day of the month} & よっか \mbox{(yok{\middel}ka)} \\\hline
    四つ & \textit{four (general counter)} & よっつ \mbox{(yot{\middel}tsu)} \\\hline
    \notyet{四月} & four + moon (month) = \textit{April} & シ{\middel}ガツ \mbox{(SHI{\middel}GATSU)} \\\hline
    \notyet{四時} & four + time = \textit{four o'clock} & よ{\middel}ジ \mbox{(yo{\middel}JI)} \\\hline
    \notyet{四十} & four + ten = \textit{forty} & よん{\middel}ジュウ \mbox{(yon{\middel}J\={U})} \\\hline
    \notyet{四百} & four + hundred = \textit{four hundred} & よん{\middel}ヒャク \mbox{(yon{\middel}HYAKU)} \\\hline
\end{somme}

\begin{sentence}
\ruby{\underline{四月}}{シガツ}に\ruby{は}{わ}\ruby{\underline{雨}}{あめ}が\ruby{\underline{多}}{おお}\underline{い}。\\\hline
    SHI{\middel}GATSU ni wa ame ga \={o}{\middel}i.\\\hline
(April) (rain) (many).\\\hline
=\textit{There's a lot of rain in April.}\\\hline
\end{sentence}
\end{multicols}
\vhrulefill{5pt}
%%--------------------------------------------------------------------
%%--------------------------------------------------------------------
\begin{multicols}{2}
\begin{glyph}{Moon (月)}{moon}\label{glyph:moon}
Moon, Month.
\end{glyph}

\begin{comps}
\comprtwocone & 月 \\\hline
\comprthreecone & Moon \\\hline
\end{comps}

\begin{remglyph}
Part of the Kanji radicals (\#{74}).\\\hline
\end{remglyph}

\begin{prons}
\pronsrtwocone & \textbf{ゲツ、ガツ (GETSU, GATSU)} & \textbf{つき (tsuki)} \\\hline
\pronsrthreecone & --- & --- \\\hline
\end{prons}

\begin{remprons}
Every \hooglig{month} everyone in the \comkun{stew kitchen} gets a Hyundai
    \comon{Getz}.\\\hline
\end{remprons}

\begin{somme}
    月 & \textit{moon; month} & つき \mbox{(tsu{\middel}ki)}\\\hline
    一月 & one + month = \textit{January} & イチ{\middel}ガツ \mbox{(ICHI{\middel}GATSU)}\\\hline
    二月 & two + month = \textit{February} & ニ{\middel}ガツ \mbox{(NI{\middel}GATSU)}\\\hline
    三月 & three + month = \textit{March} & サン{\middel}ガツ \mbox{(SAN{\middel}GATSU)}\\\hline
    四月 & four + month = \textit{April} & シ{\middel}ガツ \mbox{(SHI{\middel}GATSU)}\\\hline
    \notyet{先月} & precede + month = \textit{last month} & セン{\middel}ゲツ \mbox{(SEN{\middel}GETSU)}\\\hline
    \notyet{月見} & moon + see = \textit{moon viewing} & つき{\middel}み \mbox{(tsuki{\middel}mi)}\\\hline
\end{somme}

\begin{sentence}
\ruby{\underline{九}}{ク}\ruby{\underline{月}}{ガツ}に\ruby{\underline{月}}{つき}\ruby{\underline{見}}{み}に\ruby{\underline{行}}{い}きました。\\\hline
KU{\middel}GATSU ni tsuki{\middel}mi ni i{\middel}kimashita.\\\hline
(September) (moon viewing) (went)\\\hline
I went moon viewing in September.\\\hline
\end{sentence}
\end{multicols}
\vspace*{\fill}
%%--------------------------------------------------------------------
%%--------------------------------------------------------------------
\pagebreak
\vspace*{\fill}
\begin{multicols}{2}
\begin{glyph}{Bright (明)}{bright}\label{glyph:bright}
Bright; Clear; Light.
\end{glyph}

\begin{comps}
\comprtwocone & 日 + 月 \\\hline
\comprthreecone & Sun + Moon \\\hline
\end{comps}

\begin{remglyph}
    The \hooglig{bright} \remcom{sun} and \remcom{moon}.\\\hline
\end{remglyph}

\begin{prons}
\pronsrtwocone & \textbf{メイ (MEI)} & \textbf{あ、あか (a, aka)}\\\hline
\pronsrthreecone & ミョウ、ミン (MY\={O}, MIN) & あき (aki)\\\hline
\end{prons}

\begin{remprons}
The \hooglig{bright} \ucomon{myou} dolls, \comkun{a.k.a.} \ucomkun{akimbo} dolls, generated a \ucomon{mingled} \comon{mayhem} after the \comkun{update}.\\\hline
\end{remprons}

\begin{somme}
    明るい & \textit{bright} & あか{\middel}るい \mbox{(aka{\middel}rui)}\\\hline
    明かり & \textit{clearness} & あ{\middel}かり \mbox{(a{\middel}kari)}\\\hline
    明白 & bright + white = \textit{obvious} & メイ{\middel}ハク \mbox{(MEI{\middel}HAKU)}\\\hline
    \notyet{不明} & not + bright = \textit{unclear} & フ{\middel}メイ \mbox{(FU{\middel}MEI)}\\\hline
\end{somme}

\begin{sentence}
\ruby{\underline{月}}{つき}\ruby{は}{わ}\ruby{\underline{明}}{あか}\underline{るい}です。\\\hline
tsuki wa aka{\middel}rui desu.\\\hline
(moon) (bright)\\\hline
The moon is bright.\\\hline
\end{sentence}
\end{multicols}
\vhrulefill{5pt}
%%--------------------------------------------------------------------
%%--------------------------------------------------------------------
\begin{multicols}{2}
\begin{glyph}{Tree (木)}{tree}\label{glyph:tree}
Tree; Wood.
\end{glyph}

\begin{comps}
\comprtwocone & 木 \\\hline
\comprthreecone & Tree \\\hline
\end{comps}

\begin{remglyph}
Part of the Kanji radicals (\#{75}).\\\hline
\end{remglyph}

\begin{prons}
\pronsrtwocone & \textbf{モク、ボク (MOKU, BOKU)} & \textbf{き (ki)}\\\hline
\rye{3}{Look for モク to show up in the first position.\\\hline
	Look for ボク to show up in the second position.\\\hline
	き occurs less often, but appears with equal frequency in the first or second position.}
\pronsrthreecone & --- & こ (ko)\\\hline
\end{prons}

\begin{remprons}
From the \hooglig{tree} I was drinking \comon{mock} \comon{bock} to \comkun{keep} the \ucomkun{cobbles} of everyday life away.\\\hline
\end{remprons}

\begin{somme}
    木 & \textit{tree} & き \mbox{(ki)}\\\hline
    \notyet{木曜日} & tree + day of the week + sun (day) = \textit{Thursday} & モク{\middel}ヨウ{\middel}ビ \mbox{(MOKU{\middel}Y\={O}{\middel}BI)}\\\hline
    \notyet{木星} & tree + star = \textit{Jupiter (planet)} & モク{\middel}セイ \mbox{(MOKU{\middel}SEI)}\\\hline
    \notyet{大木} & large + tree = \textit{large tree} & タイ{\middel}ボク \mbox{(TAI{\middel}BOKU)}\\\hline
    \notyet{土木} & earth +tree = \textit{civil engineering} & ド{\middel}ボク \mbox{(DO{\middel}BOKU)}\\\hline
\end{somme}

\begin{irrr}
    木綿{\notcov} & tree + cottom = \textit{cotton (cloth)} & も{\middel}メン \mbox{(mo{\middel}MEN)}\\\hline
\rye{3}{The second kanji incidentally is not covered in this book.  Such
    characters will be included only when they form words in which the
    pronounciation of the character we are considering has an irregular reading.
    Irregular readings containing these kanji will be identified with a
    star{\notcov}. }
\end{irrr}

\begin{sentence}
\underline{あの} \ruby{\underline{高}}{たか}\underline{い} \ruby{\underline{木}}{き} の\ruby{\underline{名前}}{なまえ}が\ruby{\underline{分}}{わ}\underline{かります}か。\\\hline
ano taka{\middel}i ki no na{\middel}mae ga wa{\middel}karimasu ka.\\\hline
(that) (tall) (tree) (name) (understand).\\\hline
=\textit{Do you know the name of that tall tree?}\\\hline
\end{sentence}
\end{multicols}
\vspace*{\fill}
%%--------------------------------------------------------------------
%%--------------------------------------------------------------------
\pagebreak
\vspace*{\fill}
\begin{multicols}{2}
\begin{glyph}{Five (五)}{five}\label{glyph:five}
Five/Pent-.
\end{glyph}

\begin{comps}
\comprtwocone & 五 \\\hline
\comprthreecone & Five \\\hline
\end{comps}

\begin{remglyph}
    \textit{This is how most people feel when their work ends at five o'clock.}\\\hline
\end{remglyph}

\begin{prons}
\pronsrtwocone & \textbf{ ゴ (GO)} & ---\\\hline
\pronsrthreecone & --- & いつ (itsu)\\\hline
\end{prons}

\begin{remprons}
At \hooglig{5} o'clock we don't \ucomkun{eat soup}, we \comon{gobble}
    it.\\\hline
\end{remprons}

\begin{somme}
    五 & \textit{five} & ゴ \mbox{(GO)}\\\hline
    五人 & five + person = \textit{five people} & ゴ{\middel}ニン \mbox{(GO{\middel}NIN)}\\\hline
    五月 & five + moon (month) = \textit{May} & ゴ{\middel}ガツ \mbox{(GO{\middel}GATSU)}\\\hline
    五日 & five + sun (day) = \textit{fifth day of the month} & いつ{\middel}か \mbox{(itsu{\middel}ka)}\\\hline
    五つ & \textit{five (general counter)} & いつ{\middel}つ \mbox{(itsu{\middel}tsu)}\\\hline
    \notyet{五十} & five + ten = \textit{fifty} & ゴ{\middel}ジュウ \mbox{(GO{\middel}J\={U})}\\\hline
    \notyet{五時} & five + time = \textit{five o'clock} & ゴ{\middel}ジ \mbox{(GO{\middel}JI)}\\\hline
\end{somme}

\begin{sentence}
\ruby{\underline{五十人}}{ゴジュウニン}で\ruby{\underline{山}}{やま}え\ruby{\underline{行}}{い}\underline{きました}。\\\hline
GO{\middel}J\={U}{\middel}NIN de yama e i{\middel}kimashita.\\\hline
(fifty people) (mountain) (went)\\\hline
=\textit{Fifty (of us) went to the mountain.}\\\hline
\end{sentence}
\end{multicols}
\vhrulefill{5pt}
%%--------------------------------------------------------------------
%%--------------------------------------------------------------------
\begin{multicols}{2}
\begin{glyph}{Eye (目)}{eye}\label{glyph:eye}
Eye.\\\hline
In some instances it is read with its kun-yomi, and indicates the \textbf{``-th'' suffix} in words such as ``fourth'' and ``seventh''; it also expresses the related endings of ``first'', ``second'', ``third'', etc.
\end{glyph}

\begin{comps}
\comprtwocone & 目 \\\hline
\comprthreecone & Eye \\\hline
\end{comps}

\begin{remglyph}
Part of the Kanji radicals (\#{109}).\\\hline
\end{remglyph}

\begin{prons}
\pronsrtwocone & \textbf{モク (MOKU)} & \textbf{め (me)}\\\hline
\rye{3}{め is the more common reading of the two.}
\pronsrthreecone & ボク (BOKU) & ま (ma)\\\hline
\end{prons}

\begin{remprons}
I \hooglig{eyed} the \ucomkun{muddy} \comkun{mess} caused by the \comon{mock} \ucomon{bock}.\\\hline
\end{remprons}

\begin{somme}
    目 & \textit{eye} & め \mbox{(me)}\\\hline
    \notyet{目玉} & eye + jewel = \textit{eyeball} & め{\middel}だま \mbox{(me{\middel}dama)}\\\hline
    \notyet{八回目} & eight + rotate + eye = \textit{the eighth time} & ハチ{\middel}カイ{\middel}め \mbox{(HACHI{\middel}KAI{\middel}me)}\\\hline
    \notyet{目立つ} & eye + stand = \textit{to stand out} & め{\middel}だつ \mbox{(me{\middel}datsu)}\\\hline
    \notyet{注目} & pour + eye = \textit{notice} & チュウ{\middel}モク \mbox{(CH\={U}{\middel}MOKU)}\\\hline
\end{somme}

\begin{sentence}
\ruby{\underline{王子}}{オウジ}の\ruby{\underline{目}}{め} \ruby{は}{わ}\ruby{\underline{青い}}{あおい}です。\\\hline
\={O}{\middel}JI no me wa ao{\middel}i desu.\\\hline
(prince) (eye) (blue)\\\hline
=\textit{The prince's eyes are blue.}\\\hline
\end{sentence}
\end{multicols}
\vspace*{\fill}
%%--------------------------------------------------------------------
%%--------------------------------------------------------------------
\pagebreak
\begin{multicols}{2}
\begin{glyph}{Woman (女)}{woman}\label{glyph:woman}
Woman; Female.
\end{glyph}

\begin{comps}
\comprtwocone & 女 \\\hline
\comprthreecone & Woman \\\hline
\end{comps}

\begin{remglyph}
Part of the Kanji radicals (\#{38}).\\\hline
\end{remglyph}

\begin{prons}
\pronsrtwocone & \textbf{ジョ (JO)} & \textbf{おんな (onna)}\\\hline
\pronsrthreecone & ニョ、ニョウ (NYO, NY\={O}) & め (me)\\\hline
\end{prons}

\begin{remprons}
Many \hooglig{woman}, whose \comon{jobs} are in \ucomkun{messy} \ucomon{Neon} areas (an appropriate \ucomon{neologism}), are \comkun{on narcotic} drugs.\\\hline
\end{remprons}

\begin{somme}
    女 & \textit{woman} & おんな \mbox{(onna)}\\\hline
    \notyet{女王} & woman + king = \textit{queen} & ジョ{\middel}オウ \mbox{(JO{\middel}\={O})}\\\hline
    \notyet{少女} & few + woman = \textit{girl} & ショウ{\middel}ジョ \mbox{(SH\={O}{\middel}JO)}\\\hline
    \notyet{女心} & woman + heart = \textit{a woman's heart} & おんな{\middel}ごころ \mbox{(onna{\middel}gokoro)}\\\hline
    \notyet{雪女} & snow + heart = \textit{snow fairy} & ゆき{\middel}おんな \mbox{(yuki{\middel}onna)}\\\hline
    \notyet{男女} & man + woman = \textit{men and women} & ダン{\middel}ジョ \mbox{(DAN{\middel}JO)}\\\hline
\end{somme}

\begin{sentence}
\ruby{\underline{女}}{おんな}の\ruby{\underline{人}}{ひと}は\ruby{\underline{木}}{き}の\ruby{\underline{下}}{した}で\ruby{\underline{止}}{と}\underline{まりました}。\\\hline
onna no hito wa ki no shita de to{\middel}marimashita.\\\hline
(woman) (person) (tree) (under) (stopped).\\\hline
=\textit{The woman stopped under the tree.}\\\hline
\end{sentence}
\end{multicols}
\vhrulefill{5pt}
%%--------------------------------------------------------------------
%%--------------------------------------------------------------------
\begin{multicols}{2}
\begin{glyph}{Large (大)}{large}\label{glyph:large}
Large; Big.\\\hline
The meaning of ``important'' and ``great'' (not necessarily always in a positive sense) can also be implied in words containing this kanji.
\end{glyph}

\begin{comps}
\comprtwocone & 大 \\\hline
\comprthreecone & Large\\\hline
\end{comps}

\begin{remglyph}
 Part of the Kanji radicals (\#{37}).\\\hline
\end{remglyph}

\begin{prons}
\pronsrtwocone & \textbf{ダイ、タイ (DAI, TAI)} & \textbf{おお (\={o})}\\\hline
\rye{3}{ダイ and タイ occur with equal frequency when this kanji appears in the
	first position.  おお much less often.\\\hline When it occupies other
	positions in a compound, however, ダイ is by far the more common reading.}
\pronsrthreecone & --- & ---\\\hline
\end{prons}

\begin{remprons}
The \hooglig{large}, \comon{tight} fabric had an \comkun{awful} \comon{dye} to it.\\\hline
\end{remprons}

\begin{somme}
    大きい & \textit{large} & おお{\middel}きい \mbox{(\={o}{\middel}kii)}\\\hline
    大木 & large + tree = \textit{large tree} & タイ{\middel}ボク \mbox{(TAI{\middel}BOKU)} \\\hline
    \notyet{大好き} & large + like = \textit{to really like} & ダイ　す{\middel}き \mbox{(DAI su{\middel}ki)}\\\hline
    \notyet{大体} & large + body = \textit{in general} & ダイ{\middel}タイ \mbox{(DAI{\middel}TAI)}\\\hline
    \notyet{大雨} & large + rain = \textit{heavy rain} & おお{\middel}あめ \mbox{(\={o}{\middel}ame)}\\\hline
    \notyet{大切} & large + cut = \textit{important} & ダイ{\middel}セツ \mbox{(TAI{\middel}SETSU)}\\\hline
    \notyet{大学} & large + study = \textit{university} & ダイ{\middel}ガク \mbox{(DAI{\middel}GAKU)}\\\hline
\end{somme}

\begin{irrr}
    大人  & large + person = \textit{adult}  & おとな \mbox{(otona)}  \\\hline
\end{irrr}

\begin{sentence}
\underline{あの} \ruby{\underline{森}}{もり} にわ\ruby{\underline{大木}}{タイボク}が\ruby{\underline{多}}{おお}\underline{い}。\\\hline
    ano mori ni wa TAI{\middel}BOKU ga \={o}{\middel}i.\\\hline
(that) (forest) (large trees) (many)\\\hline
=\textit{There are many large trees in that forest}.\\\hline
\end{sentence}
\end{multicols}
%%--------------------------------------------------------------------
%%--------------------------------------------------------------------
\pagebreak
\vspace*{\fill}
\begin{multicols}{2}
\begin{glyph}{Middle (中)}{middle}\label{glyph:middle}
A versatile character that expresses the ideas of ``middle'', ``medium'', ``average'', etc.\\\hline
This kanji can also serves as an abbreviation for China.\\\hline
日 is used as the equivalent with respect to Japan.
\end{glyph}

\begin{comps}
\comprtwocone & 口 + 丨 \\\hline
\comprthreecone & Mouth + Pole (Vertical Stroke) \\\hline
\end{comps}

\begin{remglyph}
    The \hooglig{middle} \remcom{pole} will kill \remcom{dracula}.\\\hline
\end{remglyph}

\begin{prons}
\pronsrtwocone & \textbf{チュウ (CH\={U})} & \textbf{なか (naka)}\\\hline
\pronsrthreecone & ジュウ (J\={U}) & ---\\\hline
\end{prons}

\begin{remprons}
In the \hooglig{middle} of the \ucomon{juice} container is \comon{chewable}
    \comkun{knick-knacks}.\\\hline
\end{remprons}

\begin{somme}
    中 & \textit{middle} & なか \mbox{(naka)}\\\hline
    \notyet{中心} & middle + heart = \textit{center} & チュウ{\middel}シン \mbox{(CH\={U}{\middel}SHIN)}\\\hline
    \notyet{中米} & middle + rice = \textit{Central America} & チュウ{\middel}ベイ \mbox{(CH\={U}{\middel}BEI)}\\\hline
    \notyet{中国} & middle + country = \textit{China} & チュウ{\middel}ゴク \mbox{(CH\={U}{\middel}GOKU)}\\\hline
    \notyet{中立} & middle + stand = \textit{neutrality} & チュウ{\middel}リツ \mbox{(CH\={U}{\middel}RITSU)}\\\hline
    \notyet{中古} & middle + old = \textit{secondhand} & クウ{\middel}コ \mbox{(CH\={U}{\middel}KO)}\\\hline
    \notyet{水中} & water + middle = \textit{underwater/in the water} & スイ{\middel}チュウ \mbox{(SUI{\middel}CH\={U})}\\\hline
\end{somme}

\begin{sentence}
\ruby{\underline{車}}{くるま}の\ruby{\underline{中}}{なか}に\ruby{\underline{何}}{なに}が\underline{あります}か。\\\hline
kuruma no naka ni nani ga arimasu ka.\\\hline
(car) (middle) (what) (is)\\\hline
=\textit{What is in your car?}\\\hline
\end{sentence}
\end{multicols}
\vhrulefill{5pt}
%%--------------------------------------------------------------------
%%--------------------------------------------------------------------
\begin{multicols}{2}
\begin{glyph}{Eight (八)}{eight}\label{glyph:eight}
Eight.
\end{glyph}

\begin{comps}
\comprtwocone & 八 \\\hline
\comprthreecone & Eight\\\hline
\end{comps}

\begin{remglyph}
Part of the Kanji radicals (\#{12}).\\\hline
\end{remglyph}

\begin{prons}
\pronsrtwocone & \textbf{ハチ (HACHI)} & ---\\\hline
\pronsrthreecone &  & よう、やっ、や (y\={o}, ya-, ya)\\\hline
\end{prons}

\begin{remprons}
\hooglig{\comon{Hachi}} \ucomkun{yachted} in the days of \ucomkun{yore}.\\\hline
\end{remprons}

\begin{somme}
    八 & \textit{eight} & ハチ \mbox{(HACHI)}\\\hline
    八月 & eight + moon (month) = \textit{August} & ハチ{\middel}ガツ \mbox{(HACHI{\middel}GATSU)}\\\hline
    八人 & eight + person = \textit{eight people} & ハチ{\middel}ニン \mbox{(HACHI{\middel}NIN)}\\\hline
    八日 & eight + sun (day) = \textit{the eighth day of the month} & よう{\middel}か \mbox{(y\={o}{\middel}ka)}\\\hline
    八つ & \textit{eight (general counter)} & やっ{\middel}つ \mbox{(yat{\middel}tsu)}\\\hline
    \notyet{八円} & eight + circle = \textit{eight yen} & ハチ{\middel}エン \mbox{(HACHI{\middel}EN)}\\\hline
    \notyet{八時} & eight + time = \textit{eight o'clock} & ハチ{\middel}ジ \mbox{(HACHI{\middel}JI)}\\\hline
\end{somme}

\begin{sentence}
\ruby{\underline{八月}}{ハチガツ}に\ruby{\underline{秋田}}{あきた}え\ruby{\underline{行}}{い}\underline{きましょう}。\\\hline
HACHI{\middel}GATSU ni Aki{\middel}ta e i{\middel}kimash\={o}.\\\hline
(August) (Akita) (let's go)\\\hline
=\textit{Let's go to Akita in August.}\\\hline
\end{sentence}
\end{multicols}
\vspace*{\fill}
%%--------------------------------------------------------------------
%%--------------------------------------------------------------------
\pagebreak
\vspace*{\fill}
\begin{multicols}{2}
\begin{glyph}{Small (小)}{small}\label{glyph:small}
Small; Little.
\end{glyph}

\begin{comps}
\comprtwocone & 小 \\\hline
\comprthreecone & Small\\\hline
\end{comps}

\begin{remglyph}
Part of the Kanji radicals (\#{42}).\\\hline
\end{remglyph}

\begin{prons}
\pronsrtwocone & \textbf{ショウ (SH\={O})} & \textbf{こ (ko)}\\\hline
\pronsrthreecone & --- & お (o)\\\hline
\end{prons}

\begin{remprons}
    \hooglig{Small} cheese \comkun{cobbles} along the \comon{shore} will
    \ucomkun{oblige}.\\\hline
\end{remprons}

\begin{somme}
    小さい & \textit{small} & ちい{\middel}さい \mbox{(chii{\middel}sai)}\\\hline
    \notyet{小春} & small + spring = \textit{Indian Summer} & こ{\middel}はる \mbox{(ko{\middel}haru)}\\\hline
    \notyet{小国} & small + country = \textit{small country} & ショウ{\middel}コク \mbox{(SH\={O}{\middel}KOKU)}\\\hline
    \notyet{小高い} & small + tall = \textit{slightly elevated} & こ　だ{\middel}かい \mbox{(ko da{\middel}kai)}\\\hline
    \notyet{最小} & most + small = \textit{smallest} & サイ{\middel}ショウ \mbox{(SAI{\middel}SH\={O})}\\\hline
\end{somme}

\begin{sentence}
\underline{あの} \ruby{\underline{車}}{くるま} \ruby{は}{わ}\underline{とても} \ruby{\underline{小}}{ちい}\underline{さい} \underline{ですね}。\\\hline
ano kuruma wa totemo chii{\middel}sai desu ne.\\\hline
(that) (car) (really) (small) (isn't it).\\\hline
=\textit{That car is really small, isn't it?}\\\hline
\end{sentence}
\end{multicols}
\vhrulefill{5pt}
\vspace*{\fill}
%%--------------------------------------------------------------------
%%--------------------------------------------------------------------
\pagebreak
